% ============================================================
% PAGE 8-11 -- PROPOSED METHOD
% ============================================================
\section{Proposed Methodology}
\label{sec:method}

This work proposes a compiler-oriented framework in which the analysis phases
of a compiler feed a machine-learning model that makes the optimisation
decision. The framework has an \emph{offline} part (building the training
data) and an \emph{online} part (advising a new loop).

\subsection{Proposed Workflow}

The overall workflow consists of:

\begin{enumerate}
    \item Source/Input acquisition (a C function \texttt{kernel} with a loop nest).
    \item Lexical analysis (tokens).
    \item Syntax analysis (abstract syntax tree).
    \item Semantic analysis (symbol table, trip counts, strides, dependences).
    \item Intermediate representation generation (three-address code).
    \item Feature-vector construction.
    \item ML decision model $\rightarrow$ optimisation configuration.
    \item Target-code generation by GCC with the selected flags.
    \item Execution and performance evaluation.
\end{enumerate}

\subsection{Proposed Flow Diagram}

\begin{figure}[H]
\centering

\begin{tikzpicture}[
node distance=0.65cm and 0.8cm,
box/.style={
rectangle,
rounded corners,
draw,
minimum width=2.7cm,
minimum height=0.85cm,
align=center,
font=\small
},
arrow/.style={-{Stealth},thick}
]

\node[box] (src) {C loop kernel};
\node[box,right=of src] (lex) {Lexical\\analysis};
\node[box,right=of lex] (syn) {Syntax analysis\\(AST)};
\node[box,below=of syn] (sem) {Semantic and\\loop analysis};
\node[box,left=of sem] (ir) {Intermediate code\\(TAC)};
\node[box,left=of ir] (feat) {Feature\\vector};
\node[box,below=of feat] (ml) {ML decision\\model};
\node[box,right=of ml] (flags) {GCC flag\\selection};
\node[box,right=of flags] (exe) {Optimised\\executable};

\draw[arrow] (src) -- (lex);
\draw[arrow] (lex) -- (syn);
\draw[arrow] (syn) -- (sem);
\draw[arrow] (sem) -- (ir);
\draw[arrow] (ir) -- (feat);
\draw[arrow] (feat) -- (ml);
\draw[arrow] (ml) -- (flags);
\draw[arrow] (flags) -- (exe);

\end{tikzpicture}

\caption{Proposed workflow: compiler front end, feature extraction and ML-based optimisation decision.}
\label{fig:workflow}
\end{figure}

\begin{figure}[H]
\centering
\begin{tikzpicture}[
node distance=0.65cm and 0.8cm,
box/.style={rectangle,rounded corners,draw,minimum width=2.7cm,minimum height=0.85cm,align=center,font=\small},
arrow/.style={-{Stealth},thick}
]
\node[box] (gen) {Kernel generator\\(\numFamilies\ families)};
\node[box,right=of gen] (comp) {Compile with all\\\numConfigs\ configurations};
\node[box,right=of comp] (time) {Execute and time\\(5 interleaved rounds)};
\node[box,below=of time] (lab) {Safety filter and\\labelling};
\node[box,left=of lab] (fe) {Front end:\\feature extraction};
\node[box,left=of fe] (train) {Train and\\cross-validate};
\draw[arrow] (gen) -- (comp);
\draw[arrow] (comp) -- (time);
\draw[arrow] (time) -- (lab);
\draw[arrow] (lab) -- (fe);
\draw[arrow] (fe) -- (train);
\end{tikzpicture}
\caption{Offline phase: construction of the labelled data set and training.}
\label{fig:offline}
\end{figure}

\subsection{Lexical Analysis}

Lexical analysis converts a sequence of characters into tokens. The scanner
(\texttt{lexer.py}) is specified by regular expressions in the same way as a
Lex/Flex description, with token classes \texttt{KEYWORD}, \texttt{TYPE},
\texttt{IDENT}, \texttt{NUMBER}, \texttt{OP}, \texttt{PUNCT}; comments,
white space and pre-processor lines are discarded and an illegal character is
reported as a lexical error. For example, the statement

\begin{verbatim}
C[i][j] += A[i][k]*B[k][j];
\end{verbatim}

is converted into the token sequence (as printed by the implemented scanner):

\begin{verbatim}
IDENT C | [ i ] [ j ] | OP += | IDENT A | [ i ] [ k ] |
OP * | IDENT B | [ k ] [ j ] | ;
\end{verbatim}

\subsection{Syntax Analysis}

Syntax analysis determines whether the token sequence conforms to the grammar
of the language. The output is an Abstract Syntax Tree (AST). The parser
(\texttt{loop\_parser.py}) is a recursive-descent (LL(1)) parser; expressions are
parsed by precedence climbing, which corresponds to the usual left-recursion-free
expression grammar:

\begin{align*}
E &\rightarrow T\,(\,(+\,|\,-)\,T\,)^{*} &
T &\rightarrow F\,(\,(*\,|\,/)\,F\,)^{*} \\
F &\rightarrow \mathit{NUM}\;|\;\mathit{ID}\,[\,E\,]^{*}\;|\;(\,E\,)\;|\;-F &
S &\rightarrow \texttt{for}\;(\dots)\;S\;|\;\texttt{if}\;(E)\;S\;[\texttt{else}\;S]\;|\;L\;\mathit{op}\;E\,;
\end{align*}

where $\mathit{op}\in\{=,+\!=,-\!=,*\!=,/\!=\}$ and $L$ is a scalar or an array element.
The AST nodes are \texttt{For}, \texttt{If}, \texttt{Assign}, \texttt{BinOp},
\texttt{UnOp}, \texttt{Index}, \texttt{Var}, \texttt{Num} and \texttt{Call}.
Syntax errors are reported with the offending token and line.

\subsection{Semantic and Loop Analysis}

Semantic analysis checks contextual constraints: every array is declared and used
with the declared number of subscripts, every identifier is declared, and loop
bounds are constant expressions that can be evaluated from the \texttt{\#define}
constants (a simple symbol table). On the checked AST, the analyser computes the
information that drives the optimisation decision.

\paragraph{Trip counts.} For \texttt{for (i = lo; i < hi; i += s)}, the trip count is
$\lceil (hi-lo)/s\rceil$; the iteration space of a nest is the product of the trip counts.

\paragraph{Affine subscripts and strides.} Each subscript is evaluated as an
affine function $a\,v+b$ of the innermost loop variable $v$ (the function is
checked to be linear by evaluating it at three points). For an array with
dimensions $S_0,\dots,S_{n-1}$ and subscripts $e_0,\dots,e_{n-1}$ with
coefficients $a_k$ in $v$, the memory stride of the reference (in elements) is
\begin{equation}
\mathrm{stride}_v = \sum_{k=0}^{n-1} a_k \prod_{m>k} S_m .
\end{equation}
A stride of 1 is a unit-stride (vectorisable, cache-friendly) access, a stride of 0 is
loop-invariant, and a large stride (e.g.\ $S_1=M$ when the \emph{row} index is the inner loop
variable) is a cache-hostile access that loop interchange could repair.

\paragraph{Loop-carried dependences.} For a write $W$ and a read $R$ of the same
array with equal coefficients $a$ and offsets $b_W$, $b_R$ in the innermost
variable, the dependence distance is
\begin{equation}
d = \frac{b_W - b_R}{a}.
\end{equation}
If $d>0$ the read uses a value written $d$ iterations earlier (a flow dependence of
distance $d$) and the loop cannot be vectorised with a vector length larger than $d$.

\paragraph{Reductions.} A scalar (or an array element that is invariant in the inner loop)
updated with \texttt{+=}, \texttt{*=} or \texttt{x = x + \dots} is a reduction. Floating-point
reductions can only be vectorised if the compiler may re-associate the additions,
which is not allowed under strict IEEE semantics.

\subsection{Intermediate Representation}

The AST of the loop body is lowered to three-address code (\texttt{tac.py}). Each
instruction has at most one operator; simple affine subscripts remain inline.
Three-address code for the body of the matrix-multiplication loop
\texttt{C[i][j] += A[i][k]*B[k][j]} is:

\begin{align}
t_1 &= C[i][j]\\
t_2 &= A[i][k]\\
t_3 &= B[k][j]\\
t_4 &= t_2 * t_3\\
t_5 &= t_1 + t_4\\
C[i][j] &= t_5
\end{align}

The number of TAC instructions of the innermost body is used as a feature
(\texttt{tac\_len}) and the TAC of the complete loop nest, including labels and
conditional jumps, can be printed by the tool (Section~\ref{sec:results}).

\subsection{Feature Vector}

The analyser produces the \numFeatures\ features of Table~\ref{tab:features}. They
describe the size of the iteration space, the arithmetic and memory work of the body,
the access pattern, the dependence structure and the working set.

\begin{table}[H]
\centering
\caption{Static loop features used by the decision model.}
\label{tab:features}
\footnotesize
\begin{tabularx}{\textwidth}{lX}
\toprule
\textbf{Feature} & \textbf{Meaning}\\
\midrule
\texttt{dtype\_bytes} & Size of the data type (4 = float, 8 = double)\\
\texttt{nest\_depth} & Depth of the loop nest\\
\texttt{inner\_trip\_log2}, \texttt{total\_iters\_log2} & $\log_2$ of the innermost trip count and of the total iterations\\
\texttt{n\_add}, \texttt{n\_mul}, \texttt{n\_div}, \texttt{n\_flops} & Floating-point operations per innermost iteration (address arithmetic excluded)\\
\texttt{n\_loads}, \texttt{n\_stores}, \texttt{n\_arrays} & Distinct array loads, stores and arrays per iteration\\
\texttt{n\_branches} & Number of \texttt{if} statements in the body\\
\texttt{n\_unit\_stride}, \texttt{n\_nonunit\_stride}, \texttt{n\_invariant} & Number of references with stride 1, stride $\neq$ 0,1 and stride 0\\
\texttt{max\_stride\_log2} & $\log_2$ of the largest absolute stride\\
\texttt{has\_reduction} & 1 if the body contains a reduction\\
\texttt{has\_carried\_dep}, \texttt{dep\_distance} & 1 if a loop-carried flow dependence exists, and its distance\\
\texttt{footprint\_log2} & $\log_2$ of the total size in bytes of the arrays touched\\
\texttt{bytes\_per\_iter}, \texttt{arith\_intensity} & Bytes moved per iteration and flops per byte\\
\texttt{tac\_len} & Number of TAC instructions in the loop body\\
\bottomrule
\end{tabularx}
\end{table}

\subsection{Decision Space}

The model chooses one of the \numConfigs\ GCC configurations of Table~\ref{tab:configs}.
They are ordered from conservative to aggressive: \texttt{O2} is the default baseline,
the \texttt{native} variants use the instruction set of the build machine (less portable),
and \texttt{O3\_fast\_native} additionally allows floating-point re-association
(\texttt{-ffast-math}), which changes numerical results.

\begin{table}[H]
\centering
\caption{Decision space: the seven candidate optimisation configurations.}
\label{tab:configs}
\small

\begin{tabular}{ll}
\toprule
\textbf{Class} & \textbf{GCC flags} \\
\midrule
\texttt{O2} & \texttt{-O2} \\
\texttt{O3} & \texttt{-O3} \\
\texttt{O2\_unroll4} & \texttt{-O2 -funroll-loops --param max-unroll-times=4} \\
\texttt{O2\_unroll8} & \texttt{-O2 -funroll-loops --param max-unroll-times=8} \\
\texttt{O3\_native} & \texttt{-O3 -march=native} \\
\texttt{O3\_unroll\_native} & \texttt{-O3 -funroll-loops -march=native} \\
\texttt{O3\_fast\_native} & \texttt{-O3 -march=native -ffast-math} \\
\bottomrule
\end{tabular}
\end{table}


\subsection{Labelling Rule}

Every kernel is compiled and executed with all configurations. Let $T_c$ be the
measured time of configuration $c$ and $r_c$ the checksum of the program output. A
configuration is \emph{admissible} if its result agrees with the baseline:
\begin{equation}
\frac{|r_c-r_{O2}|}{|r_{O2}|} \le 10^{-3}.
\end{equation}
The label is the most conservative admissible configuration whose time is within
$\tau=5\%$ of the fastest admissible one:
\begin{equation}
y = \min_{\text{order}}\{\,c \;:\; c\ \text{admissible},\ T_c \le (1+\tau)\min_{c'\ \text{adm.}} T_{c'}\,\}.
\end{equation}
The tolerance $\tau$ keeps the label stable under measurement noise
(Section~\ref{sec:results}), and the conservative ordering prefers portable, numerically safe
flags whenever a more aggressive flag brings no real gain.

\subsection{Decision Models}

Four supervised classifiers from scikit-learn \cite{pedregosa2011} are compared: a depth-limited
\emph{Decision Tree} (interpretable, depth 5), a \emph{Random Forest} \cite{breiman2001} with
300 trees, \emph{Gradient Boosting} (100 trees of depth 3) and \emph{k-NN} on standardised features
($k=5$, distance weighting). They are compared with four baselines that do not use the features:
always \texttt{O2}, always \texttt{O3}, always \texttt{O3\_native} and the majority class.

\subsection{Proposed Algorithm}

\begin{algorithm}[H]
\caption{ML-Based Loop Optimisation Decision}
\label{alg:compiler}
\begin{algorithmic}[1]

\Require Loop kernel source $S$; trained model $f$; configurations $C$
\Ensure Compiler flags for $S$

\State $tokens \gets \textsc{Lex}(S)$
\State $ast \gets \textsc{Parse}(tokens)$
\State Check declarations, subscripts and loop bounds (semantic analysis)
\State $nest \gets$ loop nest with the largest iteration space
\State Evaluate trip counts, strides, dependence distance, reductions
\State $tac \gets \textsc{LowerToTAC}(ast)$
\State $\phi \gets$ feature vector from $nest$, $tac$ and the array declarations
\State $c^{*} \gets f(\phi)$
\State \Return \textsc{Flags}$(c^{*})$

\end{algorithmic}
\end{algorithm}

\begin{algorithm}[H]
\caption{Construction of the Training Set and Model}
\label{alg:training}
\begin{algorithmic}[1]

\Require Kernel corpus $\mathcal{K}$; configurations $C$; tolerance $\tau$
\Ensure Trained model $f$

\For{each kernel $L\in\mathcal{K}$}
    \State $\phi(L) \gets$ \textsc{ExtractFeatures}$(L)$
    \For{each configuration $c\in C$}
        \State compile $L$ with $c$
    \EndFor
    \For{round $=1$ \textbf{to} 5}
        \For{each $c\in C$}
            \State run the binary of $c$, record $T$ and the checksum
        \EndFor
    \EndFor
    \State $T_c \gets$ minimum over the rounds
    \State $y(L) \gets$ label by the rule above
\EndFor
\State Train $f$ on $\{(\phi(L),y(L))\}$ and evaluate with cross-validation
\State \Return $f$

\end{algorithmic}
\end{algorithm}

\subsection{Implementation}

The software is organised as follows (all files are provided with this report):

\begin{table}[H]
\centering
\caption{Source files of the implementation.}
\label{tab:files}
\small
\begin{tabular}{ll}
\toprule
\textbf{File} & \textbf{Role}\\
\midrule
\texttt{lexer.py} & Phase 1: scanner\\
\texttt{loop\_parser.py} & Phase 2: recursive-descent parser, AST\\
\texttt{loop\_analyzer.py} & Phase 3: semantic and loop analysis, feature extraction\\
\texttt{tac.py} & Phase 4: three-address-code generator\\
\texttt{kernel\_gen.py} & Generator of the kernel corpus and timing harness\\
\texttt{benchmark.py} & Compile/execute under all configurations\\
\texttt{build\_dataset.py} & Features, timings, safety filter and labels\\
\texttt{train\_model.py} & Training and evaluation of the models\\
\texttt{predict.py} & Command-line advisor for a new loop\\
\texttt{make\_report\_assets.py} & Generates the tables and plots of this report\\
\bottomrule
\end{tabular}
\end{table}

\newpage
